\section{Benchmark: Causal Diagnosis of Paralog Failure}

We design the benchmark as a causal measurement, not a tool ranking. The goal is to isolate and quantify the specific failure modes that prevent existing tools from aligning paralogs, so that each failure mode can be mapped to a specific algorithmic remedy.

\subsection{Experimental Setup}

\textbf{Species and gene sets.} We benchmark on four vertebrate species spanning increasing evolutionary distance from human: mouse (\textit{Mus musculus}, GRCm39, $\sim$90 Mya~\cite{waterston2002mouse}), dog (\textit{Canis familiaris}, ROS\_Cfamiliaris\_v1, $\sim$96 Mya~\cite{lindblad2005dog}), chicken (\textit{Gallus gallus}, bGalGal1, $\sim$320 Mya~\cite{hillier2004chicken}), and zebrafish (\textit{Danio rerio}, GRCz11, $\sim$435 Mya~\cite{howe2013zebrafish}). For one-to-one orthologs, we use 2,000 randomly sampled proteins per species from Ensembl Compara one-to-one orthology annotations. For paralogs, we use all proteins where the human query maps to $\geq$2 paralogous loci in the target genome, yielding $n = 488$--500 proteins per species. Ground truth is defined by Ensembl Compara orthology mapping and canonical CDS annotations per species.

\textbf{Tools evaluated.} We include six tools spanning the major algorithmic families in protein-to-genome alignment: minimizer-based spliced alignment (miniprot~\cite{li2023miniprot}, miniSplice~\cite{zhang2025minisplice}), profile-HMM alignment (Spaln3~\cite{gotoh2024spaln3}), and alignment graph methods (GTH~\cite{gremme2005genomethreader}, Scipio~\cite{keller2008scipio}, genBlastG~\cite{she2011genblastg}). We also evaluate the miniprot+miniSplice pipeline (miniprot spliced alignment refined by miniSplice's deep learning model). All tools are run with protein-level queries (no genomic region restriction) and default parameters.

\textbf{Metrics.} We evaluate using exon-level F1 score (harmonic mean of exon nucleotide precision and recall), the standard metric for spliced alignment evaluation \cite{burset1996evaluation}. We report two complementary metrics:
\begin{itemize}
    \item \textbf{Primary F1:} exon F1 at the single locus selected by the tool as its best prediction. This is the metric that matters for genome annotation---it measures whether the tool produces a correct gene model at the correct genomic location.
    \item \textbf{Best-overlap F1:} exon F1 of the tool's top-ranked hit, computed against whichever Ensembl Compara ortholog (best or alternative) it overlaps. This measures alignment quality at the locus the tool actually selected, without penalizing selection of a co-ortholog rather than the highest-identity ortholog. The gap between primary and best-overlap F1 reflects the structural divergence between the selected paralog copy and the best-identity copy.
\end{itemize}
Confidence intervals are 95\% paired bootstrap with $10^4$ resamples (seed 4242).

We distinguish between the \textbf{A-stage} (search/alignment tools that find genomic candidate regions and produce initial alignments) and the \textbf{C-stage} (spliced-alignment refinement tools that produce final exon-intron structure). The intermediate steps—chaining, splice site prediction, and exon assembly—are internal to individual tools and do not constitute a separately benchmarkable stage. The A/C naming reflects this two-part structure: A denotes the initial search/alignment step, and C denotes the final spliced-alignment refinement step. The internal selection step (B) is discussed in the error decomposition (Section~2.2) but does not correspond to an independently exchangeable tool category. The benchmark evaluates tools in their end-to-end configuration (A-stage + C-stage combined).

\subsection{The Paralog Gap}

Table~\ref{tab:benchmark} shows end-to-end results. On one-to-one orthologs, the best tools achieve 0.94--0.96 exon F1 in mammals and 0.89 in chicken---one-to-one ortholog alignment is mature. On paralogs, however, all tools show substantial degradation: the best single-tool primary F1 ranges from 0.34 (mouse) to 0.17 (zebrafish). The gap between one-to-one and paralog primary F1 is 0.55--0.80 across species. No tool exceeds 0.41 primary F1 on any species' paralog set.

The best-overlap F1 in the paralog regime (0.50--0.63) is substantially higher than primary F1---indicating that when tools select a co-ortholog locus rather than the best-identity copy, the spliced alignment at that co-ortholog can still be of reasonable quality. The primary-to-best\_overlap gap (0.07--0.24) reflects the structural divergence between the selected paralog copy and the best-identity ortholog, and already points to locus confusion (RC2).

Zebrafish is the hardest species across all tools, consistent with its greater evolutionary distance from human (mean human--zebrafish protein identity $\sim$55\% vs $\sim$80\% for human--mouse). This species gradient is diagnostic: if failure were purely a tool engineering issue, we would expect tool-specific patterns rather than a monotonic degradation with evolutionary distance.

To determine \emph{why} all tools fail on paralogs, we decompose each tool's end-to-end performance into three error categories per protein (Exp~2): whether the correct locus is found by the search stage (A), whether it is ranked first by the tool's internal selection step (B), and whether the spliced alignment meets the quality threshold given correct selection (C). Unlike A and C, which correspond to separately benchmarkable tool categories (Section~2.1), B is an internal step within each end-to-end tool---it cannot be isolated as an independent stage, but its contribution to overall failure can be quantified through the decomposition. The results are shown in Table~\ref{tab:s_exp2_abc}.

\begin{figure}[t]
    \centering
    \includegraphics[width=\columnwidth]{figures/fig1_paralog_gap.pdf}
    \caption{Primary vs.\ best-overlap exon F1 on paralogs across 6 alignment tools and TPRS. For existing tools, the primary-to-best-overlap gap (gray solid vs.\ gray hatched) quantifies locus confusion (RC2): tools can produce reasonable alignments somewhere but fail to select the correct locus. TPRS (gold) substantially closes both the primary and best-overlap gaps.}
    \label{fig:paralog_gap}
\end{figure}

\begin{table*}[t]
\centering
\caption{Complete end-to-end benchmark results (Exp~1) and TPRS performance. Mean exon F1 with 95\% paired-bootstrap CI ($10^4$ resamples, seed 4242). Paralog columns report \textbf{primary} / \emph{best-overlap} F1 (evaluated against actually-matched ortholog). $n = 2000$ (1:1), 488--500 (paralog) per species.}
\label{tab:benchmark}
\footnotesize
\setlength{\tabcolsep}{2.5pt}
\resizebox{\textwidth}{!}{%
\begin{tabular}{lcccccccc}
\toprule
Tool & \multicolumn{2}{c}{Mouse} & \multicolumn{2}{c}{Dog} & \multicolumn{2}{c}{Chicken} & \multicolumn{2}{c}{Zebrafish} \\
 & 1:1 & Paralog & 1:1 & Paralog & 1:1 & Paralog & 1:1 & Paralog \\
\midrule
\texttt{miniprot} & 0.956 & 0.326 (0.29, 0.37) / \emph{0.597 (0.55, 0.64)} & 0.944 & 0.350 (0.31, 0.39) / \emph{0.600 (0.56, 0.64)} & 0.886 & 0.400 (0.36, 0.44) / \emph{0.498 (0.46, 0.54)} & 0.867 & 0.170 (0.14, 0.20) / \emph{0.247 (0.21, 0.29)} \\
\texttt{Spaln3} & 0.763 & 0.289 (0.25, 0.33) / \emph{0.496 (0.45, 0.54)} & 0.745 & 0.268 (0.23, 0.30) / \emph{0.503 (0.46, 0.54)} & 0.713 & 0.295 (0.26, 0.33) / \emph{0.373 (0.33, 0.42)} & 0.663 & 0.095 (0.07, 0.12) / \emph{0.136 (0.11, 0.17)} \\
\texttt{GTH} & 0.827 & 0.338 (0.30, 0.38) / \emph{0.611 (0.57, 0.65)} & 0.804 & 0.344 (0.31, 0.38) / \emph{0.563 (0.52, 0.60)} & 0.776 & 0.366 (0.33, 0.40) / \emph{0.470 (0.43, 0.51)} & 0.695 & 0.135 (0.11, 0.16) / \emph{0.215 (0.18, 0.25)} \\
\texttt{Scipio} & 0.636 & 0.151 (0.12, 0.18) / \emph{0.333 (0.29, 0.37)} & 0.836 & 0.328 (0.29, 0.36) / \emph{0.489 (0.45, 0.53)} & 0.721 & 0.286 (0.25, 0.32) / \emph{0.338 (0.30, 0.37)} & 0.617 & 0.101 (0.08, 0.12) / \emph{0.134 (0.11, 0.16)} \\
\texttt{genBlastG} & 0.733 & 0.157 (0.13, 0.19) / \emph{0.300 (0.26, 0.34)} & 0.828 & 0.255 (0.22, 0.29) / \emph{0.567 (0.53, 0.60)} & 0.480 & 0.312 (0.28, 0.35) / \emph{0.443 (0.40, 0.48)} & 0.746 & 0.160 (0.13, 0.19) / \emph{0.283 (0.25, 0.32)} \\
\texttt{mini+miniSplice} & 0.959 & 0.333 (0.29, 0.37) / \emph{0.608 (0.57, 0.65)} & 0.945 & 0.351 (0.31, 0.39) / \emph{0.623 (0.58, 0.66)} & 0.892 & 0.401 (0.36, 0.44) / \emph{0.507 (0.47, 0.55)} & 0.873 & 0.173 (0.14, 0.21) / \emph{0.250 (0.21, 0.29)} \\
\midrule
\textbf{TPRS} & \textbf{0.993} & \textbf{0.793 (0.71, 0.87)} / \emph{0.836 (0.76, 0.90)} & \textbf{0.993} & \textbf{0.771 (0.69, 0.85)} / \emph{0.854 (0.79, 0.91)} & \textbf{0.994} & \textbf{0.760 (0.67, 0.84)} / \emph{0.802 (0.73, 0.87)} & \textbf{0.993} & \textbf{0.575 (0.47, 0.67)} / \emph{0.635 (0.54, 0.73)} \\
$\Delta_{\mathrm{primary}}$ vs best e2e &  & \textbf{+0.456} &  & \textbf{+0.421} &  & \textbf{+0.359} &  & \textbf{+0.401} \\
$\Delta_{\mathrm{best\text{-}overlap}}$ vs best e2e &  & \textbf{+0.225} &  & \textbf{+0.231} &  & \textbf{+0.295} &  & \textbf{+0.352} \\
\bottomrule
\end{tabular}
}
\end{table*}

\begin{table*}[!t]
\centering
\caption{Per-tool stage-level error decomposition (Exp~2). OK = exon F1 $\ge 0.95$. A = search (locus $\notin$ top-50). B = selection (in top-50, not \#1). C = alignment (\#1, F1 $< 0.95$). $\tau=0.95$, K=50. Bold $\ge 50\%$. $\Delta$A = A$_{\text{para}} - $A$_{\text{1:1}}$ (pp). 1:1 dominant: C=\texttt{miniprot},\texttt{GTH},\texttt{mini+miniS}; others A/B. Paralog: A 23/24 cells.}
\label{tab:s_exp2_abc}
\footnotesize
\setlength{\tabcolsep}{2pt}
\begin{tabular}{llrrrrrrrrr}
\toprule
Tool & Sp & \multicolumn{4}{c}{1:1 ortholog} & \multicolumn{4}{c}{Paralog (broad)} & $\Delta$A \\
\cmidrule{3-6} \cmidrule{7-10}
 & & OK\% & A\% & B\% & C\% & OK\% & A\% & B\% & C\% & (pp) \\
\midrule
\texttt{miniprot} & Mo & 82.7 & 6.9 & 0.2 & 10.2 & 26.0 & 50.0 & 1.8 & 22.2 & +43.1 \\
 & Do & 71.2 & 8.1 & 0.1 & 20.6 & 19.8 & 47.0 & 3.0 & 30.2 & +38.9 \\
 & Ch & 52.1 & 14.0 & 0.0 & 33.9 & 18.9 & 56.6 & 1.0 & 23.6 & +42.6 \\
 & Zf & 60.1 & 13.7 & 0.2 & 26.0 & \textbf{8.8} & \textbf{80.0} & 1.2 & 10.0 & +66.3 \\
\midrule
\texttt{Spaln3} & Mo & 62.6 & 27.6 & 0.0 & 9.8 & 22.8 & \textbf{61.4} & 0.0 & 15.8 & +33.8 \\
 & Do & 51.6 & 29.8 & 0.0 & 18.6 & 16.2 & \textbf{55.8} & 0.0 & 28.0 & +26.0 \\
 & Ch & 40.5 & 32.6 & 0.0 & 26.9 & 14.1 & \textbf{66.4} & 0.0 & 19.5 & +33.8 \\
 & Zf & 45.0 & 36.8 & 0.0 & 18.2 & \textbf{6.4} & \textbf{88.0} & 0.0 & 5.6 & +51.2 \\
\midrule
\texttt{GTH} & Mo & 44.5 & 6.9 & 0.2 & 48.4 & 23.4 & 50.0 & 1.8 & 24.8 & +43.1 \\
 & Do & 37.8 & 8.1 & 0.1 & \textbf{54.1} & 16.8 & 47.0 & 3.0 & 33.2 & +38.9 \\
 & Ch & 27.6 & 14.0 & 0.0 & \textbf{58.5} & 11.9 & 56.6 & 1.0 & 30.5 & +42.6 \\
 & Zf & 25.6 & 13.7 & 0.2 & \textbf{60.5} & \textbf{5.8} & \textbf{80.0} & 1.2 & 13.0 & +66.3 \\
\midrule
\texttt{Scipio} & Mo & 30.4 & \textbf{55.2} & 0.0 & 14.3 & 6.4 & \textbf{80.0} & 0.0 & 13.6 & +24.8 \\
 & Do & 26.9 & \textbf{57.2} & 0.0 & 15.9 & 7.6 & \textbf{72.4} & 0.4 & 19.6 & +15.2 \\
 & Ch & 9.3 & \textbf{64.3} & 0.0 & 26.4 & 2.9 & \textbf{80.7} & 0.0 & 16.4 & +16.4 \\
 & Zf & 5.3 & \textbf{63.5} & 0.0 & 31.1 & \textbf{0.4} & \textbf{95.8} & 0.2 & 3.6 & +32.3 \\
\midrule
\texttt{genBlastG} & Mo & 8.8 & 44.3 & 37.9 & 9.0 & 1.2 & \textbf{83.0} & 13.8 & 2.0 & +38.7 \\
 & Do & 9.6 & 37.0 & \textbf{44.9} & 8.6 & 0.8 & 48.0 & \textbf{46.4} & 4.8 & +11.0 \\
 & Ch & 2.8 & \textbf{63.6} & 24.3 & 9.2 & 1.6 & \textbf{57.4} & 36.1 & 4.9 & $-$6.2 \\
 & Zf & 7.3 & 37.4 & \textbf{42.6} & 12.7 & 1.4 & \textbf{61.8} & 33.0 & 3.8 & +24.4 \\
\midrule
\texttt{mini+miniS} & Mo & \textbf{84.0} & 6.9 & 0.2 & 8.8 & 26.6 & 50.0 & 1.8 & 21.6 & +43.1 \\
 & Do & \textbf{71.9} & 8.1 & 0.1 & 20.0 & 19.6 & 47.0 & 3.0 & 30.4 & +38.9 \\
 & Ch & \textbf{55.4} & 14.0 & 0.0 & 30.6 & 19.3 & 56.6 & 1.0 & 23.2 & +42.6 \\
 & Zf & \textbf{65.4} & 13.7 & 0.2 & 20.6 & 10.2 & \textbf{80.0} & 1.2 & 8.6 & +66.3 \\
\bottomrule
\end{tabular}

\vspace{2pt}
\footnotesize $\Delta$A range $-$6.2 to +66.3 pp; 23/24 cells $\ge +11$ pp. Spaln3--miniprot A gap (1:1): Mo +20.7, Do +21.7, Ch +18.6, Zf +23.1 pp. Robustness $\rho = 1.000$ (1:1, 9 cuts); $\rho \ge 0.935$ (para, 8/9).
\end{table*}


The stage decomposition (Table~\ref{tab:s_exp2_abc}) yields two findings. First, under the paralog regime the dominant failure mode shifts from C-stage to A-stage across all tools---for the C-dominated 1:1 tools (\texttt{miniprot}, \texttt{mini+miniSplice}, \texttt{GTH}), the error class flips from C to A, with A-stage growth ranging from $+5.6$ to $+71.5$ pp (23/24 cells $\ge +21$ pp). Second, even on 1:1 orthologs, tools exhibit heterogeneous bottlenecks: \texttt{miniprot} and \texttt{mini+miniSplice} are C-dominated, while \texttt{Spaln3} and \texttt{Scipio} are A-dominated---the same dataset stresses different tools at different stages. Both findings establish that locus-level search is the primary bottleneck in paralog alignment, and that any solution must accommodate tool-level diversity in search capability. The A/B/C decomposition identifies \emph{where} tools fail. We now isolate \emph{why,} decomposing these error patterns into two independent root causes: the combined A+B failure of locus search and selection (RC2), and the C-stage alignment ceiling imposed by cross-species sequence divergence (RC1). The forced-locus experiment (Exp~5) cleanly separates the two by removing the search problem entirely.

\subsection{Root Cause 1: Cross-Species Identity Ceiling}

To determine whether paralog failure is caused by search error (tools cannot find the correct locus) or alignment error (tools find the correct locus but cannot produce a good alignment), we conduct a forced-locus experiment (Exp~5). Each tool is run with the correct genomic locus provided as a $\pm$50~kb window around the ortholog-annotated CDS, removing the search problem entirely. If locus confusion were the only issue, paralog forced-locus F1 should recover to one-to-one forced-locus levels.

Results (Table~\ref{tab:forced_locus}) refute this hypothesis. Even with perfect locus knowledge, a human protein query against a non-human paralog locus produces substantially degraded alignments. With miniprot---the strongest C-stage tool---forced-locus paralog F1 reaches only 0.64--0.71, compared to 0.92--0.98 for forced-locus one-to-one orthologs. The 0.22--0.30 gap between these conditions is the \textbf{cross-species identity ceiling}: a fundamental limit imposed by amino acid sequence divergence between the human query protein and the target species' paralogous gene product. No C-stage alignment algorithm can overcome this because the query protein simply lacks the sequence motifs necessary to resolve exon boundaries in a diverged target. The identity gap is largest in zebrafish (0.31), consistent with its greater evolutionary distance. Notably, Scipio on zebrafish paralogs exhibits a junction-vs-exon F1 discrepancy (0.905 vs.\ 0.324, Table~\ref{tab:forced_locus}), indicating correct splice junction identification but incorrect exon boundary assembly---a pattern consistent with chaining algorithms recovering short conserved segments while misassembling the full CDS at extreme divergence.

This finding directly motivates the query-identity dimension of TPRS candidate generation: by including the target species' own protein as a query (sharing $\sim$100\% amino acid identity with the target gene), the identity barrier is eliminated.

\begin{table*}[t]
\centering
\caption{Forced-locus alignment ceiling (Exp~5, $\pm$50 kb pad). Exon F1 / junction F1 / CDS recovery. $\Delta_{\text{id}}$ = 1:1 exon F1 $-$ paralog exon F1. $n = 2000$ (1:1), 488--500 (paralog) per species. Bold = best per column.}
\label{tab:forced_locus}
\footnotesize
\setlength{\tabcolsep}{2.5pt}
\resizebox{\textwidth}{!}{%
\begin{tabular}{lcccccccccc}
\toprule
Tool & Sp & \multicolumn{3}{c}{1:1 ortholog} & \multicolumn{3}{c}{Paralog (broad)} & \multicolumn{2}{c}{Identity Gap} \\
\cmidrule{3-5} \cmidrule{6-8} \cmidrule{9-10}
 & & Exon F1 & Junc F1 & CDS & Exon F1 & Junc F1 & CDS & $\Delta$Exon & $\Delta$Junc \\
\midrule
\texttt{miniprot} & Mo & \textbf{0.976} & 0.934 & 0.973 & 0.708 & 0.658 & 0.695 & 0.268 & 0.276 \\
 & Do & \textbf{0.955} & 0.907 & 0.944 & 0.698 & 0.641 & 0.708 & 0.257 & 0.266 \\
 & Ch & \textbf{0.923} & 0.846 & 0.899 & 0.694 & 0.630 & 0.706 & 0.229 & 0.216 \\
 & Zf & \textbf{0.946} & 0.884 & 0.927 & \textbf{0.641} & 0.574 & 0.623 & 0.305 & 0.310 \\
\midrule
\texttt{Spaln3} & Mo & 0.913 & 0.815 & 0.915 & 0.719 & 0.574 & 0.720 & 0.194 & 0.241 \\
 & Do & 0.886 & 0.782 & 0.880 & 0.642 & 0.499 & 0.660 & 0.244 & 0.283 \\
 & Ch & 0.876 & 0.756 & 0.864 & 0.629 & 0.518 & 0.642 & 0.247 & 0.238 \\
 & Zf & 0.909 & 0.795 & 0.917 & 0.594 & 0.471 & 0.590 & 0.315 & 0.324 \\
\midrule
\texttt{GTH} & Mo & 0.521 & 0.424 & 0.567 & 0.337 & 0.294 & 0.411 & 0.184 & 0.130 \\
 & Do & 0.526 & 0.432 & 0.571 & 0.313 & 0.236 & 0.385 & 0.213 & 0.196 \\
 & Ch & 0.585 & 0.474 & 0.617 & 0.373 & 0.308 & 0.422 & 0.212 & 0.166 \\
 & Zf & 0.594 & 0.501 & 0.618 & 0.324 & 0.274 & 0.332 & 0.270 & 0.227 \\
\midrule
\texttt{Scipio} & Mo & 0.969 & 0.848 & 0.951 & \textbf{0.821} & 0.711 & 0.811 & 0.148 & 0.137 \\
 & Do & \textbf{0.955} & 0.863 & 0.933 & \textbf{0.807} & 0.662 & 0.819 & 0.148 & 0.201 \\
 & Ch & 0.896 & 0.738 & 0.841 & \textbf{0.850} & 0.689 & 0.856 & 0.046 & 0.049 \\
 & Zf & 0.909 & 0.731 & 0.856 & 0.324 & 0.905\textsuperscript{*} & 0.355 & 0.585 & -- \\
\midrule
\texttt{genBlastG} & Mo & 0.962 & 0.872 & 0.959 & 0.700 & 0.494 & 0.693 & 0.262 & 0.378 \\
 & Do & 0.945 & 0.859 & 0.934 & 0.664 & 0.504 & 0.675 & 0.281 & 0.355 \\
 & Ch & 0.897 & 0.740 & 0.874 & 0.610 & 0.414 & 0.604 & 0.287 & 0.326 \\
 & Zf & 0.916 & 0.744 & 0.899 & 0.537 & 0.331 & 0.496 & 0.379 & 0.413 \\
\midrule
\texttt{GeneWise} & Mo & 0.961 & 0.882 & 0.949 & 0.724 & 0.567 & 0.692 & 0.237 & 0.315 \\
 & Do & 0.938 & 0.845 & 0.919 & 0.705 & 0.557 & 0.680 & 0.233 & 0.288 \\
 & Ch & 0.865 & 0.750 & 0.813 & 0.582 & 0.451 & 0.533 & 0.283 & 0.299 \\
 & Zf & 0.889 & 0.786 & 0.839 & 0.461 & 0.319 & 0.389 & 0.428 & 0.467 \\
\midrule
\texttt{Exonerate} & Mo & 0.952 & 0.892 & 0.931 & 0.652 & 0.558 & 0.614 & 0.300 & 0.334 \\
 & Do & 0.939 & 0.877 & 0.914 & 0.666 & 0.544 & 0.643 & 0.273 & 0.333 \\
 & Ch & 0.832 & 0.747 & 0.764 & 0.550 & 0.479 & 0.508 & 0.282 & 0.268 \\
 & Zf & 0.805 & 0.736 & 0.723 & 0.431 & 0.344 & 0.360 & 0.374 & 0.392 \\
\bottomrule
\multicolumn{11}{l}{\textsuperscript{*}Scipio zebrafish paralog: Junc=0.905 vs Exon=0.324---short-exon chaining artifact inflates junction score; true alignment quality is poor. $\Delta$Junc omitted as uninterpretable.} \\
\multicolumn{11}{l}{Identity gap (miniprot): 0.229--0.305 (mean 0.265). C-stage paralog ceiling: 0.64--0.71 (miniprot).} \\
\end{tabular}
}
\end{table*}


\subsection{Root Cause 2: Locus Confusion}

With the identity ceiling quantified, we decompose the end-to-end paralog failure into two additive gaps (Fig.~\ref{fig:forced_locus}a). The \textbf{identity gap} (forced-locus one-to-one F1 $-$ forced-locus paralog F1) is the loss attributable purely to cross-species sequence divergence. The \textbf{search gap} (forced-locus paralog F1 $-$ end-to-end paralog F1) is the loss attributable to locus selection error. For miniprot, the identity gap is 0.22--0.30 and the search gap is 0.22--0.38 across species. Both gaps are substantial and of comparable magnitude, confirming that paralog alignment requires addressing \emph{both} problems simultaneously---addressing only one would leave at least 0.22 F1 on the table.

The correct-locus rate analysis (Fig.~\ref{fig:forced_locus}b, Table~S1) confirms the severity of locus confusion. In the one-to-one regime, tools correctly rank the true locus first 0.91--0.99 of the time. In the paralog regime, this rate collapses to 0.11--0.49. Tools select the wrong paralog copy more than half the time because paralogs share high sequence identity (often $>$80\% at the protein level), and raw similarity scores cannot distinguish the evolutionarily correct copy from a highly similar sibling.

The search-stage recall analysis (Exp~3) reveals a critical asymmetry: while the best A-stage tool (\texttt{miniprot}) achieves recall@50 of 0.89 on 1:1 orthologs (4-species mean), recall drops to 0.42 on paralogs---a 53\% relative decline.

\begin{table*}[!t]
\centering
\caption{A-stage search recall@50 (Exp~3). Search-OK rate = fraction of proteins with correct locus in top-50. Per-species recall@50 for both regimes. 4-sp mean K-curve shown per tool (not per-species). Bold = best per column.}
\label{tab:s_exp3_recall}
\footnotesize
\setlength{\tabcolsep}{2.5pt}
\begin{tabular}{llrrrc}
\toprule
Tool & Sp & \multicolumn{2}{c}{Recall@50} & \multicolumn{2}{c}{4-sp mean K-curve} \\
\cmidrule{3-4} \cmidrule{5-6}
 & & 1:1 & Paral. & r@1 $\rightarrow$ r@50 $\rightarrow$ r@200 & $\Delta_{1\to200}$ (type) \\
\midrule
\texttt{miniprot} & Mo & \textbf{0.931} & 0.500 & \multirow{4}{*}{0.892 $\rightarrow$ 0.893 $\rightarrow$ 0.893} & \multirow{4}{*}{+0.001 (flat)} \\
 & Do & \textbf{0.919} & \textbf{0.530} & & \\
 & Ch & \textbf{0.860} & 0.434 & & \\
 & Zf & \textbf{0.863} & 0.200 & & \\
\midrule
\texttt{mmseqs2} & Mo & 0.821 & 0.200 & \multirow{4}{*}{0.750 $\rightarrow$ 0.802 $\rightarrow$ 0.802} & \multirow{4}{*}{+0.052 (flat)} \\
 & Do & 0.814 & 0.516 & & \\
 & Ch & 0.783 & \textbf{0.465} & & \\
 & Zf & 0.790 & \textbf{0.380} & & \\
\midrule
\texttt{last} & Mo & 0.812 & 0.180 & \multirow{4}{*}{0.758 $\rightarrow$ 0.812 $\rightarrow$ 0.813} & \multirow{4}{*}{+0.055 (flat)} \\
 & Do & 0.801 & 0.518 & & \\
 & Ch & 0.759 & 0.391 & & \\
 & Zf & 0.752 & 0.334 & & \\
\midrule
\texttt{spaln3} & Mo & 0.724 & 0.386 & \multirow{4}{*}{0.683 $\rightarrow$ 0.683 $\rightarrow$ 0.683} & \multirow{4}{*}{+0.000 (flat)} \\
 & Do & 0.702 & 0.442 & & \\
 & Ch & 0.674 & 0.336 & & \\
 & Zf & 0.632 & 0.120 & & \\
\midrule
\texttt{tblastn} & Mo & 0.611 & 0.106 & \multirow{4}{*}{0.600 $\rightarrow$ 0.694 $\rightarrow$ 0.711} & \multirow{4}{*}{+0.111 (rising)} \\
 & Do & 0.706 & 0.434 & & \\
 & Ch & 0.725 & 0.355 & & \\
 & Zf & 0.736 & 0.336 & & \\
\midrule
\texttt{genBlastA}\textsuperscript{1} & Mo & 0.611 & 0.106 & \multicolumn{2}{c}{= tblastn (same HSPs)} \\
 & Do & 0.706 & 0.434 & & \\
 & Ch & 0.725 & 0.355 & & \\
 & Zf & 0.736 & 0.336 & & \\
\midrule
\texttt{genBlastG} & Mo & 0.557 & 0.170 & \multirow{4}{*}{0.170 $\rightarrow$ 0.544 $\rightarrow$ 0.605} & \multirow{4}{*}{\textbf{+0.435} (steep)} \\
 & Do & 0.630 & 0.520 & & \\
 & Ch & 0.363 & 0.426 & & \\
 & Zf & 0.626 & 0.382 & & \\
\midrule
\texttt{diamond} & Mo & 0.091 & 0.420 & \multirow{4}{*}{0.081 $\rightarrow$ 0.082 $\rightarrow$ 0.082} & \multirow{4}{*}{+0.001 (flat)} \\
 & Do & 0.074 & 0.484 & & \\
 & Ch & 0.062 & 0.129 & & \\
 & Zf & 0.040 & 0.108 & & \\
\bottomrule
\multicolumn{6}{l}{\textsuperscript{1}\texttt{genBlastA} uses the same tblastn HSPs but groups/ranks them differently (HSP-grouping vs sum-of-bitscores). Recall numbers are identical; ranking differs (see Exp~4).} \\
\multicolumn{6}{l}{\textbf{Oracle 2-tool recall@50 (paralog):} Mo \texttt{miniprot+mmseqs} 0.578; Do \texttt{last+genBlastA} 0.796; Ch \texttt{mmseqs+genBlastA} 0.680; Zf \texttt{mmseqs+genBlastA} 0.490. \textbf{4-tool union:} Mo 0.658, Do 0.812, Ch 0.709, Zf 0.526.} \\
\multicolumn{6}{l}{\textbf{Key:} \texttt{miniprot} mean recall 1:1 0.893 $\rightarrow$ paralog 0.416 (53\% drop). \texttt{genBlastG} 6$\rightarrow$3, \texttt{tblastn/genBlastA} 4$\rightarrow$6 on paralog. \texttt{genBlastG} steepest K-curve---locus in candidate pool but buried deep. Exp~2$\leftrightarrow$3 $\Delta$=0.0000.} \\
\end{tabular}
\end{table*}


No single A-stage tool exceeds 0.53 recall@50 on any species' paralog set (best: \texttt{miniprot} on dog, 0.530), and the four complementary tools selected for TPRS (\texttt{diamond}, \texttt{last}, \texttt{miniprot}, \texttt{mmseqs2}) each exhibit distinct species-level failure patterns (Table~\ref{tab:s_exp3_recall})---for instance, \texttt{diamond} achieves 0.420 on mouse but only 0.108 on zebrafish, while \texttt{mmseqs2} shows the reverse pattern (0.200 vs.\ 0.380), confirming that tool diversity is necessary for comprehensive locus coverage.

Furthermore, conditional selection analysis (Exp~4) shows that even when the correct locus is in the candidate set, tools with informative rankers (\texttt{genBlastA}) select it only 53--67\% of the time on paralogs (vs.\ 92--98\% on 1:1), and tools with structurally broken rankers (\texttt{genBlastG}) select it only 11--19\% of the time.

\begin{table*}[!t]
\centering
\caption{B-stage conditional selection (Exp~4). P = P(Select-OK $\mid$ Search-OK) at K=50 [mean rank]. $\Delta$P = 1:1 P $-$ paralog P. Type: narrow (1--2 candidates, P$\approx$1 vacuous), degraded (ranker breaks under paralog), broken (structurally broken). genBlA$\Delta$ = genBlastA-native $-$ tblastn-sumbits P @ K=10. Bold = $\Delta$P $> 0.10$.}
\label{tab:s_exp4_selection}
\footnotesize
\setlength{\tabcolsep}{2.5pt}
\begin{tabular}{llrrrlrc}
\toprule
Tool & Sp & \multicolumn{2}{c}{1:1} & \multicolumn{2}{c}{Paralog (broad)} & $\Delta$P & Type & genBlA$\Delta$ \\
\cmidrule{3-4} \cmidrule{5-6}
 & & P & Rk & P & Rk & & & \\
\midrule
\texttt{miniprot} & Mo & 0.997 & 1.00 & 0.964 & 1.04 & 0.033 & narrow & --- \\
 & Do & 0.999 & 1.00 & 0.943 & 1.06 & 0.056 & narrow & --- \\
 & Ch & 1.000 & 1.00 & 0.976 & 1.02 & 0.024 & narrow & --- \\
 & Zf & 0.997 & 1.00 & 0.940 & 1.06 & 0.057 & narrow & --- \\
\midrule
\texttt{spaln3} & Mo & 1.000 & 1.00 & 1.000 & 1.00 & 0.000 & narrow & --- \\
 & Do & 1.000 & 1.00 & 1.000 & 1.00 & 0.000 & narrow & --- \\
 & Ch & 1.000 & 1.00 & 1.000 & 1.00 & 0.000 & narrow & --- \\
 & Zf & 1.000 & 1.00 & 1.000 & 1.00 & 0.000 & narrow & --- \\
\midrule
\texttt{genBlastA} & Mo & 0.949 & 1.11 & 0.662 & 3.03 & \textbf{0.287} & degraded & +0.044 \\
 & Do & 0.980 & 1.04 & 0.658 & 2.06 & \textbf{0.322} & degraded & +0.051 \\
 & Ch & 0.958 & 1.09 & 0.668 & 2.18 & \textbf{0.290} & degraded & +0.076 \\
 & Zf & 0.923 & 1.15 & 0.531 & 3.17 & \textbf{0.392} & degraded & +0.074 \\
\midrule
\texttt{genBlastG} & Mo & 0.320 & 7.33 & 0.188 & 13.9 & 0.132 & broken & --- \\
 & Do & 0.288 & 7.59 & 0.108 & 10.6 & 0.180 & broken & --- \\
 & Ch & 0.331 & 7.22 & 0.154 & 11.4 & 0.177 & broken & --- \\
 & Zf & 0.319 & 7.72 & 0.136 & 11.6 & 0.183 & broken & --- \\
\midrule
\texttt{scipio} & Mo & 1.000 & 1.00 & 1.000 & 1.00 & 0.000 & narrow & --- \\
 & Do & 1.000 & 1.00 & 0.986 & 1.12 & 0.014 & narrow & --- \\
 & Ch & 1.000 & 1.00 & 1.000 & 1.00 & 0.000 & narrow & --- \\
 & Zf & 1.000 & 1.00 & 0.952 & 1.76 & 0.048 & narrow & --- \\
\bottomrule
\end{tabular}

\vspace{2pt}
\footnotesize
\textbf{genBlastA vs tblastn @ K=50 (paired CI):} Mo +0.068 [0.056,0.092]; Do +0.080 [0.065,0.096]; Ch +0.120 [0.097,0.143]; Zf +0.101 [0.084,0.118] (all CIs exclude 0). tblastn-maxbits never reverses sign. \textbf{Key:} \texttt{miniprot},\texttt{spaln3},\texttt{scipio} = narrow (P$\approx$1 structural). \texttt{genBlastA} = only informative ranker, degraded on paralog ($\Delta$P 0.29--0.39). \texttt{genBlastG} = broken (P $\le$ 0.33, rank $\ge$ 7). GenBlastA HSP-grouping $>$ simple combiners (C3 PASS).
\end{table*}


Exp~4 distinguishes two classes of tools. \texttt{genBlastA} and \texttt{genBlastG} produce informative candidate rankings: their per-query candidate sets contain tens of entries, P(Select$|$Search) measures genuine selection capability, and its degradation from 1:1 to paralog directly quantifies the selection component of RC2. In contrast, \texttt{miniprot}, \texttt{spaln3}, and \texttt{scipio} are \emph{narrow tools}---they output only 1--2 candidates per query, so P(Select$|$Search) $\approx$ 1 whenever the correct locus is in the candidate set. This is a structural property of small candidate set size, not evidence of strong ranking. For narrow tools, the selection stage is not separately measurable: their search success and selection success are fused into a single decision. The B-stage analysis in Exp~4 is therefore primarily informative for the ranked-output tools (\texttt{genBlastA}, \texttt{genBlastG}), while the A-stage recall analysis (Exp~3) applies to all tools equally.

Together, these results establish RC2 as a dual failure of search coverage and selection accuracy, both of which degrade sharply in the paralog regime (Tables~\ref{tab:s_exp3_recall}, \ref{tab:s_exp4_selection}).

This finding motivates two design decisions in TPRS: (1) the tool-diversity dimension---if different tools fail on different loci, combining their hits expands the candidate pool to include the correct locus for more proteins; and (2) learned joint selection---no fixed heuristic can reliably pick the correct locus, because the discriminating signal (e.g., genomic position relative to the primary ortholog) interacts non-linearly with alignment quality features.

\begin{figure}[t]
    \centering
    \includegraphics[width=\columnwidth]{figures/fig3_forced_locus_diagnosis.pdf}
    \caption{Causal diagnosis of paralog failure. (a) Forced-locus decomposition with miniprot: identity gap = 1:1 forced-locus $-$ paralog forced-locus (RC1); search gap = paralog forced-locus $-$ paralog end-to-end (RC2). Both gaps are 0.22--0.38 across species. (b) Correct-locus rate collapse from 1:1 (solid) to paralog (hatched) across all tools.}
    \label{fig:forced_locus}
\end{figure}

\subsection{Fixed Recipes Are Insufficient}

Having diagnosed two root causes, we ask: can the problem be solved by a fixed, non-learned combination of tools? We tested fixed-hybrid pipelines (Exp~6) that pair any A-stage tool from an 8-tool panel with any C-stage tool from a 7-tool panel, evaluating all 56 (A, C) combinations per cell.

\begin{table*}[t]
\centering
\caption{Stage-level oracle and fixed-hybrid pipeline evaluation (Exp~6). Best single E2E = strongest single end-to-end tool per cell. Best fixed hybrid = best (A-stage, C-stage) pair with miniprot-equivalent C-stage tools at pad=50 kb. Best fixed recipe = (A$^*$, C$^*$) achieving best fixed hybrid F1. Per-protein oracle = upper bound from choosing the optimal (A,C) pair per protein. H1 = $\Delta$(best fixed $-$ best single) with 95\% paired-bootstrap CI. H3 = $\Delta$(oracle $-$ best fixed) with CI. AC1 = H1 persists on identity $\ge 90\%$ subset; AC2 = H1 persists on $\ge 4/8$ A-tools-find-GT subset. A-panel: \texttt{miniprot}, \texttt{spaln3}, \texttt{genblastg}, \texttt{tblastn}, \texttt{diamond}, \texttt{mmseqs2}, \texttt{last}, \texttt{genblasta}. C-panel: \texttt{miniprot}, \texttt{spaln3}, \texttt{gth}, \texttt{exonerate}, \texttt{genewise}, \texttt{scipio}, \texttt{genblastg}.}
\label{tab:s_exp6_hybrid}
\footnotesize
\setlength{\tabcolsep}{2.5pt}
\begin{tabular}{lcccccccc}
\toprule
Cell & Best single & Best fixed & Best recipe & H1 $\Delta$ (CI) & H1 & Oracle & H3 $\Delta$ (CI) & H3 \\
\cmidrule{1-9}
\multicolumn{9}{l}{\footnotesize mp=miniprot, gba=genBlastA, sp3=spaln3, exo=exonerate, genbg=genBlastG} \\
\midrule
\textit{1:1 ortholog} \\
Mouse & 0.959 & 0.913 & mp+mp & $-$0.046 ($-$0.055,$-$0.037) & $\times$ & 0.952 & +0.039 (+0.032,+0.047) & $\times$ \\
Dog & 0.945 & 0.886 & mp+mp & $-$0.059 ($-$0.069,$-$0.050) & $\times$ & 0.915 & +0.029 (+0.023,+0.035) & $\times$ \\
Chicken & 0.892 & 0.804 & mp+mp & $-$0.088 ($-$0.099,$-$0.077) & $\times$ & 0.862 & +0.059 (+0.050,+0.068) & $\checkmark$ \\
Zebrafish & 0.873 & 0.823 & mp+mp & $-$0.050 ($-$0.060,$-$0.041) & $\times$ & 0.909 & +0.086 (+0.076,+0.097) & $\checkmark$ \\
\midrule
\textit{Paralog} \\
\textbf{Mouse} & 0.338 & \textbf{0.410} & mp+genbg & \textbf{+0.072} (+0.032,+0.113) & \textbf{$\checkmark$} & 0.640 & +0.230 (+0.196,+0.265) & $\checkmark$ \\
\textbf{Dog} & 0.351 & \textbf{0.504} & gba+exo & \textbf{+0.154} (+0.109,+0.200) & \textbf{$\checkmark$} & 0.733 & +0.229 (+0.198,+0.261) & $\checkmark$ \\
Chicken & 0.401 & 0.359 & gba+sp3 & $-$0.042 ($-$0.084,+0.000) & $\times$ & 0.653 & +0.294 (+0.258,+0.330) & $\checkmark$ \\
\textbf{Zebrafish} & 0.173 & \textbf{0.286} & gba+sp3 & \textbf{+0.113} (+0.075,+0.151) & \textbf{$\checkmark$} & 0.576 & +0.290 (+0.247,+0.333) & $\checkmark$ \\
\bottomrule
\end{tabular}

\vspace{4pt}
\footnotesize
\textbf{Key findings:}
\textbf{H1 (fixed-hybrid headroom):} Best fixed (A,C) recipe beats best single E2E by $\Delta = +0.07$ to $+0.15$ on 3/4 paralog cells (mouse, dog, zebrafish; chicken narrowly fails with CI upper edge at 0.00). On all 4 1:1 cells, no fixed recipe beats \texttt{minisplice\_miniprot}---the 1:1 regime is saturated.\\
\textbf{H2 (per-regime recipe divergence):} Best recipe differs between 1:1 and paralog on all 4 species. 1:1 always uses \texttt{miniprot+miniprot}; paralog switches search to \texttt{genblasta} (3/4 species) and kernel to \texttt{exonerate}/\texttt{spaln3}/\texttt{genblastg}.\\
\textbf{H3 (per-protein oracle ceiling):} Per-protein stage oracle exceeds best fixed hybrid by $\Delta = +0.23$ to $+0.29$ on all 4 paralog cells. This is the headroom that learned per-protein routing (TPRS Phase~4) is designed to capture.\\
\textbf{AC1 (identity $\ge 90\%$):} H1 persists on the high-identity subset for all 3 H1-passing cells (mouse +0.33, dog +0.30, zebrafish +0.80; n=13/69/15). Headroom is not an identity-coverage artifact.\\
\textbf{AC2 ($\ge 4/8$ A-tools find GT):} H1 persists on the easy subset for all 3 H1-passing cells (mouse +0.24, dog +0.29, zebrafish +0.33; n=98/268/170). Headroom is not just from picking a more-recall search tool.\\
\textbf{Abbreviations:} mp=\texttt{miniprot}, sp3=\texttt{spaln3}, genbg=\texttt{genBlastG}, gba=\texttt{genBlastA}, exo=\texttt{exonerate}. All best single E2E values are from Exp~1 end-to-end benchmark.
\end{table*}


Across all 8 (species, regime) cells, three findings emerge. First (H1), on 1:1 orthologs, no fixed recipe beats the best single end-to-end tool---the regime is saturated. On paralogs, the best fixed (A, C) pair beats the best single tool on 3/4 species ($\Delta = +0.07$ to $+0.15$, CIs exclude 0), but the winning recipe differs from 1:1 on every species (H2): 1:1 always uses \texttt{miniprot+miniprot}, while paralog switches search to \texttt{genBlastA} (3/4 species) and kernel to \texttt{exonerate}/\texttt{spaln3}/\texttt{genBlastG}. Second (H3), the per-protein stage oracle exceeds the best fixed hybrid by $+0.23$ to $+0.29$ on all 4 paralog cells---this is the headroom that learned per-protein routing is designed to capture, and it persists on high-identity ($\ge 90\%$) and easy ($\ge 4/8$ A-tools find GT) subsets (Table~\ref{tab:s_exp6_hybrid}, footnotes). Third, RC1 (the identity ceiling) is entirely unaddressed by tool-ensemble approaches: the best fixed recipe still uses human queries and therefore inherits the 0.22--0.30 forced-locus ceiling from Section~2.3. This establishes that addressing paralog alignment requires both learned selection (to close the H3 oracle gap) and species-matched querying (to break the identity ceiling)---motivating the two-stage design of TPRS.

\subsection{Related Work}

\paragraph{Protein-to-genome alignment tools.}
Protein-to-genome alignment has been studied for over two decades. Existing tools fall broadly into two paradigms. \textbf{Spliced-alignment methods} directly align protein sequences against genomic DNA, resolving exon--intron boundaries through dynamic programming or seed-and-extend heuristics. Early tools such as GeneWise~\cite{birney2004genewise} and Exonerate~\cite{slater2005exonerate} used full dynamic programming with heuristic speedups; later tools introduced diverse algorithmic strategies: minimizer-based chaining (miniprot~\cite{li2023miniprot}), profile-HMM spliced alignment (Spaln~\cite{gotoh2008spaln}, Spaln2~\cite{iwata2012spaln2}, Spaln3~\cite{gotoh2024spaln3}), alignment graphs (GenomeThreader~\cite{gremme2005genomethreader}, Scipio~\cite{keller2008scipio}), and BLAST-based gene model construction (genBlastG~\cite{she2011genblastg}). Deep learning has recently been applied at the spliced-alignment refinement stage: miniSplice~\cite{zhang2025minisplice,yang2025minisplice} uses a neural splice-site scorer to re-rank candidate exon chains, but its input still depends on an upstream alignment engine. \textbf{Translated-search methods}—including BLAST~\cite{altschul1990blast}, DIAMOND~\cite{buchfink2021diamond}, MMseqs2~\cite{steinegger2017mmseqs2}, and LAST~\cite{kielbasa2011last}—first translate genomic DNA in all six frames, then perform protein--protein alignment, trading splice-site precision for recall. These tools are mature and widely deployed: they serve as the protein evidence source in genome annotation pipelines (BRAKER3~\cite{gabriel2024braker3}, GALBA~\cite{bruna2023galba}, GeMoMa~\cite{keilwagen2016gemoma,keilwagen2018gemoma}) and are routinely used for cross-species gene mapping. However, their evaluation has been conducted almost exclusively on one-to-one orthologs, where sequence identity is high and the correct target locus is unambiguous. Their behavior on paralogs—where multiple genomic loci compete for the top alignment rank—has not been systematically characterized.

\paragraph{Benchmarking under data heterogeneity.}
Standard benchmarks for protein-to-genome alignment evaluate tools on aggregate sets of orthologous gene pairs without stratifying by copy number~\cite{iwata2012spaln2,scalzitti2020g3po,freedman2025annotation,hou2024searchassessment}. While this yields a useful aggregate metric, it masks systematic failure modes: when one-to-one orthologs dominate numerically, their high scores conceal near-zero performance on multi-copy genes. Iwata and Gotoh's Spaln2 evaluation~\cite{iwata2012spaln2} observed that tool rankings can flip across species—a phenomenon our causal decomposition explains—but did not isolate paralog-specific error budgets. The principle of stratified evaluation under data heterogeneity is well established in other fields: BAliBASE~\cite{thompson2005balibase} introduced stratification by sequence identity and structural features for multiple sequence alignment; CASP~\cite{moult2018casp13} separates template-based from free-modeling targets in protein structure prediction. Our stratified benchmark applies this principle to protein-to-genome alignment for the first time, decomposing failure into independent causal components.

\paragraph{Learned selection and annotation-derived priors.}
The use of machine learning for tool or model selection has been explored in adjacent domains. FrugalGPT~\cite{chen2023frugalgpt} introduced cascade-tier routing among large language models per query, using the outcomes of cheaper models to condition the selection of expensive ones. TPRS transfers this principle to protein-to-genome alignment: cheap A-stage features (tool agreement, genomic position) condition the selection of a single (locus, query-type) pair whose C-stage alignment is returned. In bioinformatics, Aliro~\cite{lacava2023aliro} uses automated machine learning to recommend algorithms at the dataset level, whereas TPRS makes per-protein decisions using query-specific features. The XGBoost-driven feature importance in TPRS (Table~S3) provides interpretability absent from most learned selectors. On the annotation side, Ensembl Compara~\cite{vilella2009ensembl,herrero2016ensembl} has long curated orthology and paralogy relationships, while TOGA~\cite{kirilenko2023toga} uses an XGBoost model to classify aligned regions as ortholog, paralog, or processed pseudogene. TPRS differs from post-hoc classifiers in that it performs learned selection \emph{before} producing the final spliced alignment, integrating annotation-derived priors (species-matched proteins, curated orthology) as inputs to the decision rather than as labels applied afterward.

\paragraph{Positioning of this work.}
Three gaps emerge from this landscape. First, no prior work has systematically evaluated protein-to-genome alignment tools on paralogs with stratified, per-category metrics and causal decomposition of failure into independent root causes. Second, existing alignment methods are either entirely annotation-free—and therefore subject to the cross-species identity ceiling identified in our forced-locus diagnosis—or embedded in full annotation pipelines that are not evaluated on per-protein alignment metrics. Third, learned selection at the per-protein, per-locus level has not been applied to the protein-to-genome alignment problem. TPRS occupies a distinct point in this space: an alignment method, evaluated on alignment metrics, whose use of annotation-derived priors is motivated by causal diagnosis of \emph{where and why} annotation-free approaches fail—not by a preference for using annotations. The species-matched protein directly addresses the identity ceiling (RC1), while learned routing over an expanded candidate pool addresses locus confusion (RC2). Every architectural decision is testable through ablation (Section~4.2), distinguishing TPRS from approaches where annotation priors are introduced as post-hoc heuristics without causal justification.

\begin{figure*}[!t]
    \centering
    \includegraphics[width=\textwidth]{figures/fig2_tprs_pipeline.pdf}
    \caption{TPRS: an alignment method with annotation-derived priors. Stage 1 (Candidate Generation) expands the search space along three dimensions---tool diversity (Phase~1), query identity (Phase~3), and prior knowledge (Phase~2)---to build a candidate pool of (locus, query-type) pairs. Stage 2 (Learned Joint Selection) extracts 25 features per candidate, scores them with an XGBoost regressor, and selects the pair with the highest predicted exon F1 per protein.}
    \label{fig:tprs_pipeline}
\end{figure*}
